% Partie : théorie de la démonstration
% Chapitre : déduction naturelle
% Section : régularité et substitution

\documentclass[../../../include/open-logic-section]{subfiles}

\begin{document}

\olfileid[fr]{pt}{nat}{qua}

\olsection{Preuves régulières et substitution}

Les règles de quantification demandent quelques précautions en raison
des conditions de constante propre, aussi appelées conditions
d'\emph{eigenvariable}, de \Elim{\lexists} et de~\Intro{\lforall}.
On peut résoudre la plupart des difficultés en veillant à ne jamais
réutiliser la !!{constant}~$c$ de ces inférences. Donnons une définition
précise.

\begin{defn}
Une preuve de déduction naturelle~$\delta$ est \emph{régulière} si :
\begin{enumerate}
\item aucune constante propre~$a$ ne sert de constante propre à
plus d'une inférence \Elim{\lexists} ou~\Intro{\lforall} ;
\item si $a$ est la constante propre d'une inférence
\Elim{\lexists} ou~\Intro{\lforall}, elle n'apparaît que dans la
sous-preuve immédiate menant, respectivement, à la prémisse mineure
ou à l'unique prémisse.
\end{enumerate}
\end{defn}

\begin{lem}\ollabel{lem:subst-regular}
Supposons que $\delta$ soit régulière et se termine par~$!C$,
que $c$ soit une !!{constant} qui ne sert pas de constante propre
dans~$\delta$, et que $t$ soit un terme clos ne contenant aucune
constante propre de~$\delta$. Alors $\Subst{\delta}{t}{c}$,
obtenue en remplaçant toutes les occurrences de~$c$ par~$t$
dans~$\delta$, est une preuve régulière. Les substitutions
se font sans capture, conformément aux conventions du calcul.
Sa formule finale est $\Subst{!C}{t}{c}$.
Si $!A^x$ est une hypothèse ouverte de~$\delta$,
$\Subst{!A}{t}{c}^x$ est une hypothèse ouverte de
$\Subst{\delta}{t}{c}$.
\end{lem}

\begin{proof}
Par récurrence sur $\pheight{\delta}$.
Si $\pheight{\delta} = 0$, la preuve consiste en une seule
hypothèse~$!A^x$ ; $\Subst{\delta}{t}{c}$ est alors
$\Subst{!A}{t}{c}^x$.

Si $\pheight{\delta} > 0$, nous distinguons les cas selon
la dernière inférence de~$\delta$. Les règles propositionnelles
se traitent directement et sont laissées en exercice.
Considérons les règles de quantification.
\begin{enumerate}
\item Si la dernière inférence est $\Intro{\lforall}$,
$\delta$ a la forme
\[
    \AxiomC{}
    \RightLabel{$\delta'$}
    \DeduceC{$!A(a,c)$}
    \RightLabel{\Intro{\lforall}}
    \UnaryInfC{$\lforall[x][!A(x,c)]$}
    \DisplayProof
    \]
L'hypothèse de récurrence donne une preuve régulière
$\Subst{\delta'}{t}{c}$ de $!A(a,t)$.
Comme $a$ est une constante propre de~$\delta$, elle n'apparaît
pas dans~$t$. La dernière inférence de
$\Subst{\delta}{t}{c}$,
\[
    \AxiomC{}
    \RightLabel{$\Subst{\delta'}{t}{c}$}
    \DeduceC{$!A(a,t)$}
    \RightLabel{\Intro{\lforall}}
    \UnaryInfC{$\lforall[x][!A(x,t)]$}
    \DisplayProof
    \]
est donc correcte : la substitution n'introduit~$a$ ni dans
la conclusion ni dans les hypothèses ouvertes, où elle était
déjà absente par la condition de constante propre.

\item Si la dernière inférence est $\Elim{\lforall}$,
$\delta$ a la forme
\[
    \AxiomC{}
    \RightLabel{$\delta'$}
    \DeduceC{$\lforall[x][!A(x,c)]$}
    \RightLabel{\Elim{\lforall}}
    \UnaryInfC{$!A(s(c),c)$}
    \DisplayProof
    \]
L'hypothèse de récurrence donne une preuve régulière
$\Subst{\delta'}{t}{c}$ de $\lforall[x][!A(x,t)]$.
Le terme~$s(c)$ de la conclusion peut contenir~$c$, mais
ce n'est pas obligatoire. La dernière inférence de
$\Subst{\delta}{t}{c}$,
\[
    \AxiomC{}
    \RightLabel{$\Subst{\delta'}{t}{c}$}
    \DeduceC{$\lforall[x][!A(x,t)]$}
    \RightLabel{\Elim{\lforall}}
    \UnaryInfC{$!A(s(t),t)$}
    \DisplayProof
    \]
est correcte, et
$!A(s(t),t) \ident \Subst{!A(s(c),c)}{t}{c}$.
Le cas de \Intro{\lexists} est analogue.

\item Si la dernière inférence est $\Elim{\lexists}$,
$\delta$ a la forme
\[
    \AxiomC{}
    \RightLabel{$\delta_1'$}
    \DeduceC{$\lexists[x][!A(x,c)]$}
    \AxiomC{$\Discharge{!A(a,c)}{x}$}
    \RightLabel{$\delta_2'$}
    \DeduceC{$!C$}
    \DischargeRule{\Elim{\lexists}}{x}
    \BinaryInfC{$!C$}
    \DisplayProof
    \]
L'hypothèse de récurrence fournit les deux preuves régulières
$\Subst{\delta_1'}{t}{c}$ et $\Subst{\delta_2'}{t}{c}$,
avec substitution dans leurs conclusions et leurs hypothèses.
La première conclut à la formule existentielle substituée ;
la seconde conclut à la substitution de~$!C$, sous l'hypothèse
ouverte~$!A(a,t)$.

Puisque $a$ est une constante propre de~$\delta$, elle n'apparaît
pas dans~$t$. La dernière inférence de
$\Subst{\delta}{t}{c}$ est donc
\[
    \AxiomC{}
    \RightLabel{$\Subst{\delta_1'}{t}{c}$}
    \DeduceC{$\lexists[x][!A(x,t)]$}
    \AxiomC{$\Discharge{!A(a,t)}{x}$}
    \RightLabel{$\Subst{\delta_2'}{t}{c}$}
    \DeduceC{$\Subst{!C}{t}{c}$}
    \DischargeRule{\Elim{\lexists}}{x}
    \BinaryInfC{$\Subst{!C}{t}{c}$}
    \DisplayProof
    \]
et elle est correcte :
$\Subst{!A(x,t)}{a}{x} \ident \Subst{!A(a,c)}{t}{c}$.
La constante~$a$ n'apparaît ni dans la conclusion substituée,
ni dans la prémisse majeure, ni dans les hypothèses ouvertes
autres que la classe déchargée, car elle était absente
des formules correspondantes avant substitution et n'est
pas introduite par~$t$.
\end{enumerate}
\end{proof}

\begin{prop}\ollabel{prop:regular}
Si $\delta$ est une preuve dans \Log{N1c} ou \Log{N1i},
il existe une preuve régulière~$\delta'$ de même structure,
donc de même hauteur, qui ne diffère de $\delta$
que par le remplacement de constantes propres.
\end{prop}

\begin{proof}
Pour cette preuve seulement, disons qu'une inférence à constante
propre de~$\delta$ est \emph{propre} si sa constante n'est
la constante propre d'aucune autre inférence et n'apparaît
ailleurs dans~$\delta$ que dans la sous-preuve immédiate
autorisée par la définition de régularité.
Dans le cas contraire, disons qu'elle est \emph{sale}.

Soit $n$ le nombre d'inférences sales de~$\delta$.
Nous raisonnons par récurrence sur~$n$.
Si $n=0$, toutes les inférences à constante propre sont propres :
$\delta$ est régulière et il n'y a rien à modifier.
Sinon, choisissons une inférence \emph{sale} la plus haute possible,
par exemple une \Intro{\lforall}.
La sous-preuve~$\delta_1$ qui se termine par cette inférence
a la forme
\[
    \AxiomC{}
    \RightLabel{$\delta_2$}
    \DeduceC{$!A(c)$}
    \RightLabel{\Intro{\lforall}}
    \UnaryInfC{$\lforall[x][!A(x)]$}
    \DisplayProof
    \]
La condition de constante propre garantit que $c$ n'apparaît
ni dans~$\lforall[x][!A(x)]$ ni dans une hypothèse ouverte.
La preuve $\delta_2$ est régulière : toutes ses inférences
à constante propre sont propres, puisque l'inférence sale
choisie est la plus haute. De plus, $c$ ne sert pas de
constante propre dans~$\delta_2$ : une telle réutilisation
rendrait sale l'inférence correspondante.


Soit $a$ une !!{constant} n'apparaissant pas dans~$\delta$.
Par le \olref{lem:subst-regular}, $\Subst{\delta_2}{a}{c}$
est une preuve régulière de $!A(a)$.
Comme aucune hypothèse ouverte de~$\delta_2$ ne contient~$c$,
aucune hypothèse ouverte de~$\Subst{\delta_2}{a}{c}$
ne contient~$a$. Nous pouvons donc appliquer~$\Intro{\lforall}$.
Remplaçons $\delta_1$ dans~$\delta$ par la preuve $\delta_1'$ :
\[
    \AxiomC{}
    \RightLabel{$\Subst{\delta_2}{a}{c}$}
    \DeduceC{$!A(a)$}
    \RightLabel{\Intro{\lforall}}
    \UnaryInfC{$\lforall[x][!A(x)]$}
    \DisplayProof
    \]
L'inférence sale choisie est devenue propre ; seul le paramètre
$c$ a été remplacé par~$a$ dans sa sous-preuve.
Ce renommage ne rend sale aucune inférence auparavant propre.
La preuve résultante possède donc au plus $n-1$ inférences sales,
et l'hypothèse de récurrence s'applique.
\end{proof}

\begin{prob}
Décrire le cas de la preuve de la \olref{prop:regular}
où l'inférence sale la plus haute est~\Elim{\lexists}.
\end{prob}

Nous n'invoquerons plus explicitement cette proposition :
nous supposerons les preuves considérées régulières.
Au besoin, leurs constantes propres peuvent être renommées
pour les rendre régulières.

Le \olref{lem:subst-regular} et la \olref{prop:regular} valent aussi
pour \Log{N2c} et~\Log{N2i}. Les preuves sont laissées en exercice.














\end{document}
